\slajdid{08-s037}
\begin{frame}[label=08-s037]
U svrhu poređenja različitih logičkih kola uvešćemo pojam srednjeg kašnjenja $t_p$
\[t_p=\frac{t_{pLH}+t_{pHL}}2\]
Jasno je da ćemo se sa stanovišta inženjerstva truditi da vremena $t_{pHL}$ i $t_{pLH}$ budu što je manja moguća i približno jednaka. Primena kaže, pošto ne znamo kakva će se promena dešavati da izlazni signal moramo da „sačekamo“ uzimajući u obzir kašnjenje
\[t_p=\max(t_{pLH},t_{pHL})\]
Uočiti da nejednakost kašnjenja u izlaznom signalu izaziva i promene trajanja impulsa i pauze.
\end{frame}
